\documentclass[a4paper, 14pt]{extarticle}
\usepackage{gost}
\begin{document}
\presect{Ход работы}
\section{Система управления обучением (Moodle)}
Данные системы были созданы для дистанционного обучения. Такие
платформы созданы для того, чтобы учителя могли свободно настраивать
задания, оценивать и даже выставлять оценки. Самой известной на
сегодняшний день является платформа Moodle. \par
{\bfseries{Moodle}} --- это абсолютно бесплатная система управления
курсами с открытым исходным кодом, предназначенная для
эффективного онлайн-обучения.
Благодаря потребности в образовании и открытому исходному коду
(\href{https://github.com/moodle/moodle}{GitHub}), существует
множество сайтов, которые открыты по всему миру, в одной только России
открыто 4343 сайта (\href{https://stats.moodle.org}{Moodle
statistics}). Основным языком, на котором она написана, является PHP.
\href{https://moodle.org}{Official Moodle site}
\subsection{Назначение системы}
Данные CMS нужны в первую очередь для учебных заведений, которые
имеют множество студентов, а также для дистанционного обучения.
Бывает проблема, что если студентов очень много или обучение
предполагает дистанционный вариант, то такие платформы, как Moodle
--- целое спасение. Преподаватель может выставить задание или тест со
своими сроками сдачи. Также преподаватель может проверять задания
на платформе и выставить за неё оценку (Некоторые такие сайты могут
иметь свой "электронный" дневник). На сайте можно публиковать
объявления, учебную информацию и записи лекций.
\subsection{Описание}
\subsubsection{Назначение системы, решаемые задачи, целевая аудитория.}
Многие такие сайты (Moodle) имеют монолитную архитектуру. Существует
множество плагинов, которые позволяют расширить возможности,
соответственно, существуют и плагины, которые могут добавлять API.
\subsubsection{Архитектура}
В основном CMS писались (и пишутся) на языке PHP. Это скриптовый язык
с динамическими типами, который появился ещё в 1994 году
(\href{https://www.php.net/manual/en/history.php.php}{PHP History}) и
по сей день является актуальным языком для веб-разработки. Из-за
развитых инструментов PHP для управления базами данных (PDO), Moodle
может работать со множеством баз данных: PostgreSQL, MariaDB, MySQL,
Oracle Database, Microsoft SQL Server.
\subsubsection{Функциональные подсистемы и их взаимосвязь.}
Moodle может использовать такие подсистемы, как Web Services API
(REST, SOAP, XML-RPC). Существуют ещё такие службы для интеграции
контента на сайт: LTI (главный стандарт для подключения внешних
образовательных инструментов (антиплагиат)), SCORM (протокол для
передачи прогресса и результатов прохождения тестов), xAPI (Стандарт
для сбора подробной аналитики об учебном опыте).
\subsection{Модель данных}
Moodle может содержать до 200 таблиц. В них входят таблицы с
конфигурацией проекта, таблицы с пользователями, ролями, статистикой,
таблицы для блоков (UI), групп пользователей и вопросов.
По умолчанию Moodle содержит следующие таблицы:
\begin{itemize}
  \item Configuration
    \begin{itemize}
      \item config
        %                \begin{itemize}
        %                    \item "id" BIGINT NOT NULL,
        %                    \item "name" NVARCHAR(255) NOT NULL,
        %                    \item "value" NVARCHAR(-1) NOT NULL
        %                \end{itemize}
      \item config\_log
        %                \begin{itemize}
        %                    \item "id" BIGINT NOT NULL,
        %                    \item "userid" BIGINT NOT NULL,
        %                    \item "timemodified" BIGINT NOT NULL,
        %                    \item "plugin" NVARCHAR(100) NULL,
        %                    \item "name" NVARCHAR(100) NOT NULL,
        %                    \item "value" NVARCHAR(-1) NULL,
        %                    \item "oldvalue" NVARCHAR(-1) NULL
        %                \end{itemize}
      \item config\_plugins
        %                \begin{itemize}
        %                    \item "id" BIGINT NOT NULL,
        %                    \item "plugin" NVARCHAR(100) NOT NULL,
        %                    \item "name" NVARCHAR(100) NOT NULL,
        %                    \item "value" NVARCHAR(-1) NOT NULL
        %                \end{itemize}
    \end{itemize}

  \item User and Profiles
    \begin{itemize}
      \item user
        %                \begin{itemize}
        %                    \item "id" BIGINT NOT NULL,
        %                    \item "auth" NVARCHAR(20) NOT NULL,
        %                    \item "confirmed" SMALLINT NOT NULL,
        %                    \item "policyagreed" SMALLINT NOT NULL,
        %                    \item "deleted" SMALLINT NOT NULL,
        %                    \item "suspended" SMALLINT NOT NULL,
        %                    \item "mnethostid" BIGINT NOT NULL,
        %                    \item "username" NVARCHAR(100) NOT NULL,
        %                    \item "password" NVARCHAR(32) NOT NULL,
        %                    \item "idnumber" NVARCHAR(255) NOT NULL,
        %                    \item "firstname" NVARCHAR(100) NOT NULL,
        %                    \item "lastname" NVARCHAR(100) NOT NULL,
        %                    \item "email" NVARCHAR(100) NOT NULL,
        %                    \item "emailstop" SMALLINT NOT NULL,
        %                    \item "icq" NVARCHAR(15) NOT NULL,
        %                    \item "skype" NVARCHAR(50) NOT NULL,
        %                    \item "yahoo" NVARCHAR(50) NOT NULL,
        %                    \item "aim" NVARCHAR(50) NOT NULL,
        %                    \item "msn" NVARCHAR(50) NOT NULL,
        %                    \item "phone1" NVARCHAR(20) NOT NULL,
        %                    \item "phone2" NVARCHAR(20) NOT NULL,
        %                    \item "institution" NVARCHAR(40) NOT NULL,
        %                    \item "department" NVARCHAR(30) NOT NULL,
        %                    \item "address" NVARCHAR(70) NOT NULL,
        %                    \item "city" NVARCHAR(120) NOT NULL,
        %                    \item "country" NVARCHAR(2) NOT NULL,
        %                    \item "lang" NVARCHAR(30) NOT NULL,
        %                    \item "theme" NVARCHAR(50) NOT NULL,
        %                    \item "timezone" NVARCHAR(100) NOT NULL,
        %                    \item "firstaccess" BIGINT NOT NULL,
        %                    \item "lastaccess" BIGINT NOT NULL,
        %                    \item "lastlogin" BIGINT NOT NULL,
        %                    \item "currentlogin" BIGINT NOT NULL,
        %                    \item "lastip" NVARCHAR(45) NOT NULL,
        %                    \item "secret" NVARCHAR(15) NOT NULL,
        %                    \item "picture" BIGINT NOT NULL,
        %                    \item "url" NVARCHAR(255) NOT NULL,
        %                    \item "description" NVARCHAR(-1) NULL,
        %                    \item "descriptionformat" SMALLINT NOT NULL,
        %                    \item "mailformat" SMALLINT NOT NULL,
        %                    \item "maildigest" SMALLINT NOT NULL,
        %                    \item "maildisplay" SMALLINT NOT NULL,
        %                    \item "htmleditor" SMALLINT NOT NULL,
        %                    \item "autosubscribe" SMALLINT NOT NULL,
        %                    \item "trackforums" SMALLINT NOT NULL,
        %                    \item "timecreated" BIGINT NOT NULL,
        %                    \item "timemodified" BIGINT NOT NULL,
        %                    \item "trustbitmask" BIGINT NOT NULL,
        %                    \item "imagealt" NVARCHAR(255) NULL,
        %                    \item "screenreader" SMALLINT NOT NULL
        %                \end{itemize}
      \item user\_enrolments
        %                \begin{itemize}
        %                    \item "id" BIGINT NOT NULL,
        %                    \item "status" BIGINT NOT NULL,
        %                    \item "enrolid" BIGINT NOT NULL,
        %                    \item "userid" BIGINT NOT NULL,
        %                    \item "timestart" BIGINT NOT NULL,
        %                    \item "timeend" BIGINT NOT NULL,
        %                    \item "modifierid" BIGINT NOT NULL,
        %                    \item "timecreated" BIGINT NOT NULL,
        %                    \item "timemodified" BIGINT NOT NULL
        %                \end{itemize}
      \item user\_info\_category
        %                \begin{itemize}
        %                    \item "id" BIGINT NOT NULL,
        %                    \item "name" NVARCHAR(255) NOT NULL,
        %                    \item "sortorder" BIGINT NOT NULL
        %                \end{itemize}
      \item user\_info\_data
        %                \begin{itemize}
        %                    \item "id" BIGINT NOT NULL,
        %                    \item "userid" BIGINT NOT NULL,
        %                    \item "fieldid" BIGINT NOT NULL,
        %                    \item "data" NVARCHAR(-1) NOT NULL,
        %                    \item "dataformat" SMALLINT NOT NULL
        %                \end{itemize}
      \item user\_info\_field
        %                \begin{itemize}
        %                    \item "id" BIGINT NOT NULL,
        %                    \item "shortname" NVARCHAR(255) NOT NULL,
        %                    \item "name" NVARCHAR(-1) NOT NULL,
        %                    \item "datatype" NVARCHAR(255) NOT NULL,
        %                    \item "description" NVARCHAR(-1) NULL,
        %                    \item "descriptionformat" SMALLINT NOT NULL,
        %                    \item "categoryid" BIGINT NOT NULL,
        %                    \item "sortorder" BIGINT NOT NULL,
        %                    \item "required" SMALLINT NOT NULL,
        %                    \item "locked" SMALLINT NOT NULL,
        %                    \item "visible" SMALLINT NOT NULL,
        %                    \item "forceunique" SMALLINT NOT NULL,
        %                    \item "signup" SMALLINT NOT NULL,
        %                    \item "defaultdata" NVARCHAR(-1) NULL,
        %                    \item "defaultdataformat" SMALLINT NOT NULL,
        %                    \item "param1" NVARCHAR(-1) NULL,
        %                    \item "param2" NVARCHAR(-1) NULL,
        %                    \item "param3" NVARCHAR(-1) NULL,
        %                    \item "param4" NVARCHAR(-1) NULL,
        %                    \item "param5" NVARCHAR(-1) NULL
        %                \end{itemize}
      \item user\_lastaccess
        %                \begin{itemize}
        %                    \item "id" BIGINT NOT NULL,
        %                    \item "userid" BIGINT NOT NULL,
        %                    \item "courseid" BIGINT NOT NULL,
        %                    \item "timeaccess" BIGINT NOT NULL
        %                \end{itemize}
      \item user\_preference
        %                 \begin{itemize}
        %                     \item "id" BIGINT NOT NULL,
        %                     \item "userid" BIGINT NOT NULL,
        %                     \item "name" NVARCHAR(255) NOT NULL,
        %                     \item "value" NVARCHAR(1333) NOT NULL
        %                 \end{itemize}
      \item user\_private\_key
        %                 \begin{itemize}
        %                     \item "id" BIGINT NOT NULL,
        %                     \item "script" NVARCHAR(128) NOT NULL,
        %                     \item "value" NVARCHAR(128) NOT NULL,
        %                     \item "userid" BIGINT NOT NULL,
        %                     \item "instance" BIGINT NULL,
        %                     \item "iprestriction" NVARCHAR(255) NULL,
        %                     \item "validuntil" BIGINT NULL,
        %                     \item "timecreated" BIGINT NULL
        %                 \end{itemize}
    \end{itemize}
    \image{img/moodle-user.png}{Пользователи Moodle}{pic:MoodleUser}
  \item Роли пользователей
    \begin{itemize}
      \item role
      \item capabilities
      \item context
      \item role\_allow\_assign
      \item role\_allow\_override
      \item role\_allow\_switch
      \item role\_assignments
      \item role\_capabilities
      \item role\_context\_levels
      \item role\_names
      \item role\_sortorder
    \end{itemize}
    \image{img/moodle-role.png}{Роли пользователей}{pic:UserRole}
  \item Группы пользователей
    \begin{itemize}
      \item groups
      \item groups\_member
      \item groupings
      \item groupings\_groups
    \end{itemize}
    \image{img/moodle-group.png}{Группы пользователей}{pic:UserGroup}
\end{itemize}
\subsection{Механизмы обеспечения надёжности}
Moodle имеет несколько методов управления доступом.
\begin{itemize}
  \item Ролевая модель --- позволяет гибко настроить права для
    пользователя (студент, преподаватель, гость).
  \item Методы аутентификации --- сайт на Moodle может включать в себя
    несколько способов входа (через почту Google, VK и др.), а также
    поддерживает другие протоколы аутентификации (OAuth).
  \item Защита от XSS и SQL-инъекций
  \item Безопасность загрузки и просмотра файлов --- При открытии
    файла идёт обращение к скрипту autoload.php в который
    GET-параметром передаётся путь до файла и уже этот скрипт решает,
    как отдать файл. Если это файл кода, то он его просто даст на
    скачивание, если же это .pdf документ или же .docx (если сайт
    поддерживает открытие документов такого формата), то откроется в браузере.
\end{itemize}

\subsection{Оценка масштабируемости и производительности}
При высокой нагрузке сайт Moodle может проседать в
производительности из-за своей монолитности. Есть плагины, которые
могут предоставить REST API и с помощью него разделить ответственность
(использовать сайт Moodle как backend и написать свой фронтенд), но в
этом есть минус в том, что разработка такого решения займёт много
времени (Moodle очень много что предоставляет из коробки, в случае же
с фронтендом, писать и реализовывать придётся крайне много). \par
Сайт отлично внедряется в учебный процесс. В техникуме, в школе или в
вузе --- неважно. Почти все сайты учебных заведений сделаны на Moodle.
Помимо учебных сайтов, существуют также сайты и с научными работами.

\subsubsection{Особенности интеграции с другими системами.}
Moodle предоставляет несколько уровней интеграции. Прежде всего, это
встроенные веб-сервисы (Web Services): REST (основной и рекомендуемый
вариант), а также устаревшие SOAP и XML-RPC. Администратор включает
веб-сервисы и выпускает токен доступа; внешнее приложение обращается
к endpoint webservice/\allowbreak rest/\allowbreak server.php,
передавая wstoken, имя функции
и параметры, а ответ получает в формате JSON или XML. Через эти
функции можно программно создавать пользователей и курсы, управлять
зачислением, получать оценки и списки курсов. Тот же API использует
официальное мобильное приложение Moodle.

Для связи с другими образовательными сервисами применяются открытые
стандарты. Протокол LTI (особенно LTI 1.3 с расширением LTI
Advantage) позволяет подключать внешние инструменты — системы
проверки на плагиат, видеоконференции, интерактивные задания — с
передачей оценок и списков участников прямо в журнал Moodle. Формат
SCORM 1.2/2004 даёт возможность импортировать готовые учебные модули
ZIP-архивами и отслеживать прогресс, затраченное время и результаты
прохождения. Стандарт xAPI (Tin Can) в связке с внешним хранилищем
LRS (Learning Record Store) собирает подробную аналитику учебных
действий, в том числе вне самой платформы. Для обмена данными также
используются IMS QTI и Moodle XML (банк вопросов), CSV (пользователи
и оценки), подписки календаря iCal и ленты RSS.

Аутентификация может быть вынесена во внешнюю систему: поддерживаются
LDAP, SAML 2.0, Shibboleth, CAS, OAuth 2/OpenID Connect и подключение
к внешней базе данных, что позволяет организовать единый вход (SSO) и
централизованное управление учётными записями. Система плагинов и
событийная модель (Event API) дают возможность реагировать на
события, запускать фоновые задачи и передавать данные в электронный
дневник, систему аналитики или корпоративные сервисы. К ограничениям
относятся монолитность ядра, зависимость плагинов от версии Moodle,
необходимость использовать HTTPS и надёжно хранить токены, а также
сложности при массовой синхронизации через API — в таких случаях
интеграционную логику обычно выносят в отдельный сервис-посредник и
кэшируют часто запрашиваемые данные.
\newpage
\presect{Заключение}
\subsection{Сильные и слабые стороны}
\subsubsection{Сильные стороны}
\begin{enumerate}
  \item Открытый исходный код
  \item Высокая гибкость в настройке
  \item Богатый академический функционал
  \item Мультиязычность
  \item Сбор статистики
\end{enumerate}
\subsubsection{Слабые стороны}
\begin{enumerate}
  \item Затраты на плагины (Не все плагины, которые можно скачать в
    Moodle бесплатные)
  \item Сложный интерфейс (перегружен)
  \item Монолитная архитектура (и плюс и минус)
  \item Почти полное отсутствие поддержки
\end{enumerate}

\subsection{Контрольные вопросы}
\begin{enumerate}[label=\textbf{\arabic*.}, itemsep=1.5em, leftmargin=*]
  \item \textbf{Дайте определение информационной системы и
    перечислите её основные компоненты.}

    Информационная система (ИС) --- это связанный комплекс данных,
    технологических средств и процессов, который собирает, хранит,
    обрабатывает и выдаёт информацию для решения определённых задач.

    Основные компоненты ИС:
    \begin{itemize}[label={---}, leftmargin=1.5em, itemsep=0.5em]
      \item \textbf{Аппаратное обеспечение} (серверы, компьютеры,
        сетевое оборудование);
      \item \textbf{Программное обеспечение} (ОС, СУБД, прикладные программы);
      \item \textbf{Базы данных} (организованные совокупности
        структурированной информации);
      \item \textbf{Персонал} (пользователи, администраторы, разработчики).
    \end{itemize}

  \item \textbf{Чем облачная архитектура ИС отличается от
    традиционной клиент-серверной.}

    В традиционной архитектуре ресурсы развёртываются на физических
    серверах организации, тогда как облачная архитектура предполагает
    использование виртуализированных ресурсов, предоставляемых
    сторонним провайдером по требованию, с оплатой по факту использования.
  \item \textbf{Какие механизмы используются для обеспечения
    отказоустойчивости корпоративных ИС.}

    \begin{itemize}
      \item Резервирование
      \item Кластеризация (объединение нескольких серверов в единую
          логическую группу, где при отказе одного сервера нагрузка
        распределяется на другой).
      \item Балансировка нагрузки (метод схож с прошлым, но нагрузка
        распределяется равномерно между серверами)
      \item Микросервисная архитектура
      \item Логирование
    \end{itemize}

  \item \textbf{Что такое API, и почему он важен для интеграции систем?}

    API --- это «язык общения» между различными сервисами для их
    интеграции. Нет единого API, их много. Если
    говорить про веб-разработку, то самым известным является REST API.
    Сутью его является отдача специальной структуры с данными,
    которая называется JSON. Клиент отправляет запрос с
    определённым типом (POST, GET, PUT, PATCH, DELETE), сервер
    обрабатывает его и отдаёт ответ с кодом ответа (и может быть
    также с JSON). Если же говорить про прикладную разработку, то тут
    сложнее. Многие программы с открытым программным кодом (ffmpeg)
    могут использоваться нативно (просто вызвав его через терминал
    (командную строку)) или предоставлять некоторые API (определённые
    библиотеки), которые могут быть использованы уже в приложении
    (известная библиотека ffmpeg для Python). Ещё хорошим
    примером будет Win32, который предоставляет разработчикам
    возможность использовать некоторые функции Windows в своей
    программе (создать или удалить директорию в Windows до стандарта
    C++17). Если же брать более глубоко, то API может быть и
    таким, как
    \href{https://learn.microsoft.com/en-en/cpp/build/x64-calling-convention?view=msvc-170}{Соглашение
      о
    вызовах в Windows x64}. Это соглашение вызовов и возвратов для
    языка ассемблера.

  \item \textbf{В чём преимущества и недостатки микросервисной
    архитектуры перед монолитной?}

    Преимущества:
    \begin{itemize}[label={---}, leftmargin=1.5em, itemsep=0.5em]
      \item Распределение мощностей

        Каждый сервис выполняет свою индивидуальную задачу.
      \item Отказоустойчивость

        Если один сервис выйдет из строя, остальные будут продолжать работать

      \item Дешевизна
      \item Простота в обслуживании (не всегда)
    \end{itemize}
    Недостатки:
    \begin{itemize}[label={---}, leftmargin=1.5em, itemsep=0.5em]
      \item Сложность развёртки
      \item Сложность в обслуживании
      \item Дороговизна
    \end{itemize}

  \item \textbf{Какие показатели характеризуют эффективность
    информационной системы?}

    \begin{itemize}
      \item Количество уникальных пользователей
      \item Прибыль
      \item Быстродействие системы
      \item Низкие затраты
    \end{itemize}

  \item \textbf{Какими методами можно оценить безопасность
    веб-ориентированной ИС?}

    \begin{itemize}
      \item Метод статического анализа исходного кода (метод белого ящика)
      \item Метод динамического тестирования (метод чёрного ящика)
      \item Пентест (моделирование кибератак)
      \item Фаззинг (подача заранее некорректных данных ради
        выявления неожиданного поведения).
    \end{itemize}
\end{enumerate}

\end{document}
